\documentclass[leqno]{jsarticle}
\usepackage{amssymb}
\usepackage{amsthm}
\usepackage{amsmath}
\usepackage{comment}
	%可換図式用
		%\usepackage[all]{xy}
	%重くなるので使わいときはコメント化しておく。
%定理環境 thmはあまり使わずpropをなるべく使う。
%主張は基本的に証明の中で使い、そのため番号を付けない。
%注意環境を付けてみた。
\newtheorem{thm}{定理}
\newtheorem{prop}[thm]{命題}
\newtheorem{cor}[thm]{系}
\newtheorem{lem}[thm]{補題}
\newtheorem{df}[thm]{定義}
\newtheorem{exam}[thm]{例}
\newtheorem{fact}[thm]{事実}
\newtheorem*{claim}{主張}
\newtheorem*{cau}{注意}
\renewcommand{\proofname}{{\bf 証明}}
%矢印たち。
%\toは\rightarrowと同じ。\getsは\leftarrowと同じ。同値は\iff
%上向き矢はuparrow逆はdownarrow
\newcommand{\To}{\Longrightarrow}
\newcommand{\Gets}{\Longleftarrow}
%黒板太字
\newcommand{\nn}{\mathbb{N}}
\newcommand{\zz}{\mathbb{Z}}
\newcommand{\qq}{\mathbb{Q}}
\newcommand{\rr}{\mathbb{R}}
\newcommand{\cc}{\mathbb{C}}
%無限大
\newcommand{\mgn}{\infty}
%ギリシャン
\newcommand{\dpst}{\displaystyle}
\newcommand{\ep}{\varepsilon}
\newcommand{\al}{\alpha}
\newcommand{\be}{\beta}
\newcommand{\de}{\delta}
\newcommand{\la}{\lambda}
\newcommand{\La}{\Lambda}
\newcommand{\ga}{\gamma}
\newcommand{\lil}{\la\in \La}
%商集合
\newcommand{\qt}[2]{#1/\! #2}
%enumerate環境
\renewcommand{\labelenumi}{{\rm (\arabic{enumi})}}
%以降alignじゃなくてenumerate を使え。
%なんか演算
\newcommand{\card}{\text{{\rm card}}}
\newcommand{\supp}{\text{{\rm supp}}}
%なんか演算２
\newcommand{\bd}{\text{{\rm Bdry}}}
\newcommand{\cl}{\text{{\rm CL}}}
\newcommand{\nai}{\text{{\rm Int}}}
\newcommand{\m}{\mathfrak{m}}
%差集合は\setminusで統一する。等号の否定は\neq
%真部分集合は\subsetneqq　偏微分は\partial
%ディスプレイスタイル。\dfracでディスプレイ分数
%空集合は\Oを使う！！！！！！！！！！！！

\title{Noble's theorem}
\author{一色三良(concious77)}
\date{\today}
			\begin{document}
\maketitle
この文章ではNobleの定理と呼ばれるコンパクトハウスドルフ空間の一つの特徴づけを紹介する。事実ばっかり出てくると思うが、カントール集合のあたりの話以外はすべて電波通信の記事の中で参照できる。

\begin{fact}[tube lemma]\label{1}\footnote{「tube lemma」参照}
$X$を位相空間、$Y$の部分集合$A$をコンパクトとする。更に$x\in X$と$X\times Y$の開集合$O$が
\[
\{x\}\times A\subset O
\]
を充たしているとする。このとき$X$における$x$の開近傍$N$と$Y$における$A$の近傍$M$が存在して
\[
\{x\}\times A\subset N\times M\subset O
\]
を充たす。
\end{fact}
\begin{df}連続関数$f:X\to Y$が次の条件$(\star)$充たすとき$f$を完全写像と呼ぶ\footnote{完全写像に全射性を仮定する流儀もある。}
\begin{description}
\item[($\star$)]任意の$y\in Y$について$f^{-1}(y)$は$X$のコンパクト集合かつ$f$は閉写像
\end{description}
\end{df}

次の事実の証明は簡単であると思う。
\begin{fact}\label{2}
コンパクトハウスドルフ空間の間の連続写像$f:X\to Y$について、$f$は閉集合であり、かつ任意の$y\in Y$について$f^{-1}(y)$はコンパクトである。つまり$f$は完全写像である。
\end{fact}

次の命題は実は完全写像の特徴づけにもなっているのだが、この文章では以下に述べる程度のことしか使わない。
\begin{prop}\label{3}
$X,Y$を位相空間、$f:X\to Y$を完全写像、$Z$を位相空間とする。このとき、
\[
1_Z\times f:Z\times X\to Z\times Y
\]
は閉写像となる。\footnote{実際はより強く完全写像となる}
\end{prop}
\begin{proof}
$A\subset Z\times X$を閉写像とし、$(a,b)\not\in (1_Z\times f)(A)$とする。このとき、
\[
(1_Z\times f)^{-1}(a,b)=\{a\}\times(f^{-1}(b))
\]
となり、更に明らかに
\[
A\cap((1_Z\times f)^{-1}(a,b))=\O
\]
となる。つまり
\[
(1_Z\times f)^{-1}(a,b)\subset (Z\times X)\setminus A
\]
である。$(Z\times X)\setminus A$が開集合であることから、Tube lemmaより、$a$の開近傍$N$と$f^{-1}(b)$の開近傍$M$が存在して
\[
(1_Z\times f)^{-1}(a,b)=\{a\}\times(f^{-1}(b))\subset N\times M\subset (Z\times X)\setminus A
\]
を充たす。ここで、$V=f_!(M)=Y\setminus f(X\setminus M)$と置くと、$f$が閉写像であることと、$f^{-1}(b)\subset M$から、$V$は$M$の開近傍である。\footnote{「閉写像云々」参照}さて
\[
y\in f_!(M)\iff f^{-1}(y)\subset M
\]
が成り立つことなどから\footnote{「閉写像云々」参照}
\[
(N\times V)\cap((1_Z\times f)(A))=\O
\]
が分かる。$N\times V$は$(a,b)$の開近傍であることから、結局$(1_Z\times f)(A)$が閉集合であることが分かる。\par
よって$1_Z\times f$は閉写像である。
\end{proof}

\begin{df}\label{4}
$2=\{0,1\}$に離散位相を入れて、直積空間$D=2^{\aleph_0}$を考える。この空間をカントール空間と呼ぶ。この空間はミドルサードの議論で構成されるあのカントール集合と同相である。\footnote{このことを見るには、カントール集合の元が3進展開したときどのような形になるのかということを思い出せばよい。}
\end{df}

\begin{fact}\label{5}
$D$から単位区間$I=[0,1]$への全射連続写像が存在する。\footnote{このことを見るには、2進展開を考えればよい。}
\end{fact}

\begin{fact}\label{6}
$X$をコンパクトハウスドルフ空間とし、その位相的重みが$w(X)\le\mathfrak{m}$とする。このとき$X$は$I^{\mathfrak{m}}$に埋め込まれる。\footnote{「完全正則空間とか」参照}
\end{fact}
\begin{cau}
$w(X)$で位相空間の重みを表す。つまり位相空間$X$の開基の最小濃度である。
\end{cau}

\begin{thm}\label{7}
位相空間$X$と任意の無限基数$\mathfrak{m}$について以下は同値
\begin{eqnarray}
&X\times 2^{\m}は正規\\ 
&X\times D^{\m}は正規\\ 
&X\times I^{\m}は正規\\ 
&w(Y)\le \m となる任意のコンパクトハウスドルフ空間KについてX\times K は正規
\end{eqnarray}
\end{thm}
\begin{proof}
$\aleph_0\times \m=\m$から$2^{\m}\approx D^{\m}$となるので、$(1)\iff(2)$はわかる。\par
$(2)\To(3)$は全射$D\to I$が存在することと、上の命題\ref{3}から全射$X\times D^{\m}\to X\times I^{\m}$が閉写像となり、正規空間の閉像が正規になることからわかる。\par
$(3)\To (4)$は$w(Y)\le \m$となるコンパクトハウスドルフ空間は$I^{\m}$に埋め込めることと正規空間の閉集合は正規であることからわかる。\par
$(4)\To(1)$は自明。
\end{proof}


\begin{fact}[玉野の定理]\label{tam}\footnote{「パラコンパクト性など：アドベントカレンダー2015」参照}
$T_{3.5}$空間$X$について以下は同値
\begin{eqnarray}
&&Xはパラコンパクト\\
&&任意のコンパクトハウスドルフ空間KについてX\times KはT_4
\end{eqnarray}
\end{fact}


\begin{fact}[Stone]\label{stone}\footnote{「非可算個の可算離散空間の直積について」参照}
非空な$T_1$位相空間の族$\{X_{\lambda}\}_{\lambda\in \Lambda}$に於いて$\displaystyle \prod_{\lambda\in \Lambda}X_{\lambda}$が正規\footnote{このとき自動的に$T_4$にもなる。}なら$\Lambda$の高々可算な元を除いて$X_{\lambda}$は可算コンパクトである。
\end{fact}

\begin{fact}[Arens-Dugundji]\label{AD}\footnote{「可算コンパクトとメタコンパクト」参照}
可算コンパクトかつメタコンパクトな空間はコンパクトである。
\end{fact}
\begin{thm}[Noble]
位相空間$X$は任意の基数$\m$について$X^{\m}$が$T_4$であるならば、$X$はコンパクトハウスドルフである。
\end{thm}
\begin{proof}
$X$が空集合、または一点のみからなる空間のときは、定理の主張は自明であるので、以下$\card(X)\ge 2$と仮定する。\par
さて、$\m=1$として$X$が$T_4$であることが分かる。\par
$X^{\aleph_1}$も$T_4$なので上の事実\ref{stone}から$X$は可算コンパクトであることが分かる。\par
$X\times X^{\m}$も$T_4$なので、これの閉部分空間である$X\times 2^{\m}$も$T_4$であるから$\m$の任意性と事実\ref{6}、定理\ref{7}と事実\ref{tam}から$X$はパラコンパクトであることが分かる。そして事実から$X$はコンパクトであることが分かる。\footnote{$X$はパラコンパクトなので特にメタコンパクトである。ところで事実\ref{AD}の証明の仮定をメタコンパクトの代わりにパラコンパクトハウスドルフにまで強めると事実\ref{AD}の証明はより簡単になる。}
\end{proof}














































	
\begin{thebibliography}{9}
\bibitem{2}児玉之行,永見啓応,位相空間論,岩波書店,1974
\end{thebibliography}




			\end{document}



\begin{comment}
[可換図式]
\[\xymatrix{
&   & (2) \ar@{=>}[rd]&  &  &  & (5) \ar@{=>}[rd]&\\ 
& (1)\ar@{=>}[ru] \ar@{=>}[rd]&  &(4)\ar@{=>}[r]&(7)& (1)\ar@{=>}[ru] \ar@{=>}[rd]& & (7)&\\ 
&   & (3) \ar@{=>}[ur]&  &  &  &(6) \ar@{=>}[ur]&
}\]

[\bigin \endを使わない方法]

{\prop[aa] aaa}
で
\begin{prop}[aa]
aaa
\end{prop}
と同じ\df \thm \proof でも同様

[直和]
$\amalg$

[同値関係でよく使う波線]
\sim

[引数付きの新命令]
\newcommand{\prn}[1]{\|#1\|_p}

[番号のリセット]

\setcounter{equation}{0}


[字体の変更]

{\rm hogehoge}と使う。

\rm ローマン
\it イタリック
\sl 斜体

[supp]

\newcommand{\supp}{\text{{\rm supp}}}

[中央揃えの改行]

	\begin{eqnarray*}
		hoge1&=&hogehoge
		hoge2&=&hogehofe

	\end{eqnarray*}

&&の部分で揃えられる。


[\tag]
\begin{equation}
f(x) = ax^2 + bx + c \tag{1.2.3}
\end{equation}



[場合分け]

	\begin{cases}
		hoge1 & \text{状況１}\\
		hoge2 & \text{状況２}\\
		・
		・
		・
	\end{cases}



\end{comment}





